\documentclass{pgnotes}

\title{Document databases}

\begin{document}

\maketitle

\section*{Required reading}

\begin{enumerate}
\item \citet[Chapter 3]{perkins:2018:seven}
\end{enumerate}

\section{Document database characteristics}

Within a single \textbf{instance}, a \textbf{database} may have 1 or more \textbf{collections} ($\approx$ tables) each with 1 or more \textbf{documents} ($\approx$ rows).


We will focus on MongoDB:
\begin{itemize}
\item Can be installed locally and as-a-service on the cloud.
\item Query language: JavaScript
\item Conceptually stores documents as JSON, internally BSON.
\item The document is the unit of storage and recall.
\end{itemize}


\section{Basic operations}

Connection using \texttt{mongo} client program.
\begin{minted}{bash}
# use mongo <dbname> to connect to <dbname> database
mongo blog 
mongo shop 
\end{minted}
Usable via Python / other code using driver.

Once connected, database is in scope as \texttt{db}.

\subsection{Creation}

We can insert a document into a collection, creating the collection if it does not already exist:

\begin{minted}{javascript}
# insert the document into the products collection
db.products.insert({
    "name": "Wooden train set",
    "description": "Wooden train set painted in the Irish Rail colours.",
    "price": 10000
})
\end{minted}

Nested data possible within document:

\begin{minted}{javascript}
db.products.insert({
    "name": "Wooden train set",
    "description": "Wooden train set painted up to match operator colours.",
    "price": 10000,
    "variants": [
        "Irish",
        "British Rail",
        "Trenitalia"
    ]
});

\end{minted}

\subsection{List all documents in collection}

We can use the collection's \texttt{find()} method to get all documents in the collection.
\begin{minted}{javascript}
db.products.find()
\end{minted}

\subsection{Finding specific object by ID}

When each object is created, it receives a system-generated ID.
\begin{minted}{javascript}
db.product.find({"_id":ObjectId("6087fb3e7f26283da1384e45")});
\end{minted}

\subsection{Number of objects in collection}

\begin{minted}{javascript}
db.product.count();
\end{minted}

\subsection{Finding object by field values}

Can extend the find-by-id idea to get any other field.
\begin{minted}{javascript}
db.product.find({"price":500})
\end{minted}

\subsection{Remove a document}

Removing a \textit{document} is almost identical to the find operation.

\begin{minted}{javascript}
// remove a single document
db.product.remove({"_id":ObjectId("6087fb3e7f26283da1384e45")});

// remove all that match a query
db.product.remove({"price": 500});
\end{minted}

\section{Python access}

MongoDB uses the PyMongo driver that maps well to python structures.
Only basic usage is shown here.
Consult other sources for detailed examples. 

\subsection{Connection}

PyMongo driver provides Connection:
\begin{minted}{python}
# driver
from pymongo import MongoClient

# connect to dbname database
client = MongoClient("mongodb://localhost/dbname")
\end{minted}

\subsection{Insert}

Insert uses the standard python dictionary and list structures.

\begin{minted}{python}
# python dictionary with values
new_post = {
    "title":"Bring back dialup",
    "author": "Peadar Grant",
    "themes": ["history","internet"]
}

# insert
db.posts.insert_one(new_post)
\end{minted}

\subsection{Finding}

Finding follows the same idiom as direct connection.
JSON List / dictionary structures are transparently converted to their Python equivalents.

\subsection{Remove}

Removal follows same pattern as JavaScript. 

\section{Replication}

\section{Sharding}

\bibliographystyle{plainnat}
\bibliography{bibliography}

\end{document}